\section{Step 80 Implement Paeth with its specified ties}
Paeth chooses one neighbor; it does not return the estimate itself. Fill \texttt{row\_paeth} using left, upper and upper-left. Form the estimate left plus upper minus upper-left, calculate its absolute distance from each neighbor, and select the smallest distance. Break equal distances in the order left, upper, upper-left.

Use signed arithmetic for these small differences. Their magnitude is bounded by the byte-valued neighbors, so a negative intermediate is ordinary predictor work. Reducing a difference to a byte before taking its absolute value would change its meaning. An equivalent formulation compares the absolute values of upper minus upper-left, left minus upper-left, and left plus upper minus twice upper-left.

Place the chosen neighbor in \texttt{\$t4} and converge on \texttt{row\_paeth\_ready}; do not route this observation through the Sub region. On the first byte of the worked Paeth row, upper is 25, left and upper-left are zero, and the selected predictor is 25. Residual five will restore 30. Include tie examples when checking this region; visually similar images do not establish the required tie order. The scalar path is now complete: the worked call returns zero and 35, with every sample matching the earlier row table. The temporary whole-program caller still stops before packing.

\paragraph{Observe} Use \texttt{s80-1}; stop before \texttt{row\_paeth\_ready}. Inspect the state described above before continuing.

\Needspace{7\baselineskip}
\paragraph{Check} What does Paeth return?
\begin{itemize}
  \item[A.] The selected neighboring byte.
  \item[B.] The estimate regardless of its distance.
  \item[C.] The sum of all three neighbors.
\end{itemize}
